% !TeX encoding = UTF-8
% !TeX spellcheck = pl_PL

\def\filename{rozdzial2}

\chapter{Wpływ rozwoju sztucznej inteligencji na medycynę}

Dynamiczny rozwój technologii w XXI wieku wraz z rosnącym zapotrzebowaniem pacjentów na efektywne rozwiązania w sektorze opieki zdrowotnej, istotnie wpłynął na transformację współczesnej medycyny w ostatnich dekadach \cite{CHAPMAN20153}.
Innowacje takie jak elektroniczna dokumentacja medyczna, telemedycyna oraz technologie przenośne znacznie usprawniły dostępność i wygodę świadczeń usług zdrowotnych, a przełomowe rozwiązania nanotechnologii i druku 3D zrewolucjonizowały dziedzinę implantologii i inżynierii tkankowej \cite{HALEEM202370, mahara_revolutionising_2023, thacharodi_revolutionizing_2024}.

Szczególne znaczenie w współczesnej medycynie zyskały technologie oparte na sieciach neuronowych i sztucznej inteligencji, ponieważ umożliwiają one przetwarzanie i analizę dużych wolumenów danych klinicznych, często w czasie rzeczywistym.
Narzędzia wykorzystujące sztuczną inteligencję w medycynie można podzielić na dwie kategorie - systemy wirtualne oraz fizyczne.
Do wirtualnych zastosowań sztucznej inteligencji zalicza się narzędzia monitorujące zdrowie i analizujące dane, natomiast fizyczne zastosowania AI w medycynie obejmują roboty wspierające lub wykonujące rutynowe procedury chirurgiczne, a także inteligentne protezy sterowane adaptacyjnie \cite{amisha_overview_2019, sezer_application_2020}.

Wykorzystanie technologii opartych na sztucznej inteligencji jest spójne z paradygmatem medycyny P4, czyli medycyny personalizowanej, predyktywnej, prewencyjnej i partycypacyjnej, której celem jest umożliwienie kompleksowego podejścia do tematu zdrowia poprzez analizę rozległych zbiorów danych.
W praktyce oznacza to odejście od reaktywnego modelu opieki zdrowotnej na rzecz modelu proaktywnego, skoncentrowanego na profilaktyce oraz wczesnym wykrywaniu potencjalnych chorób.
Stan zdrowia pacjenta w medycynie 4P jest opisywany poprzez spersonalizowaną chmurę szczegółowych danych, takich jak pomiary biochemiczne, ocenę narażenia na czynniki ryzyka (np. stres, zanieczyszczenia) oraz dane genetyczne.
Inteligentne technologie medyczne wspomagają autonomię oraz wygodę pacjenta jednocześnie zmniejszając biurokrację i koszty usług medycznych poprzez ograniczenie dokumentacji automatyzację procesów \cite{glos_medycyna_2017, alonso_predictive_2019, briganti_artificial_2020}.

W ciągu ostatniej dekady algorytmy oparte na sztucznej inteligencji w postaci przenośnych urządzeń znalazły szerokie zastosowanie w kardiologii poprzez obserwowanie i analizę parametrów fizjologicznych oraz rejestrację sygnału elektrokardiograficznego.
Zdalne monitorowanie zapisu EKG z wykorzystaniem takich urządzeń jak Fitbit czy Apple Watch 4 jak i aplikacji mobilnych, np. Kardia, umożliwia wczesne wykrywanie migotania przedsionków oraz wspiera zarówno pacjentów oraz lekarzy w dokładniejszej ocenie stanu zdrowia \cite{amisha_overview_2019, riganti_artificial_2020, halcox_assessment_2017}.
Systemy przenośnych urządzeń monitorujących zsynchronizowanych z smartfonem i opartych na analizie AI sprawdziły się również w przypadku endokrynologii, dostarczając stały przepływ informacji na temat poziomu glukozy we krwi pacjentów z cukrzycą, oraz w neurologii w celu wczesnego wykrywania ataków epilepsji \cite{briganti_artificial_2020}.

Modele sztucznej inteligencji znajdują szerokie zastosowanie w diagnostyce, obejmującej między innymi takie techniki jak radiografia rentgenowska, rezonans magnetyczny, endoskopia, ultrasonografia oraz patologia cyfrowa.
Podejście to określane jest mianem diagnostyki komputerowo wspomaganej - CADx, czyli \textit{computer-aided diagnosis}.
W pulmonologii sztuczna inteligencja jest wykorzystywana do interpretacji wyników badań czynnościowych płuc.
Oprogramowanie implementujące AI, zwłaszcza z ekspertyzą pulmonologów, charakteryzuje się wyższą dokładnością diagnostyczną niż samodzielna analiza specjalistów \cite{das_collaboration_2023, topalovic_artificial_2019}.
Konwolucyjne sieci neuronowe (CNN), czyli modele głębokiego uczenia specjalizujące się w analizie obrazów, sprawdziły się w przypadku wykrywania zapalenia płuc, raka płuc, wirusa COVID-19 oraz innych anomalii na zdjęciach rentgenowskich klatki piersiowej \cite{rajaraman_visualization_2018, 10.1117/12.221630, ismael_deep_2021, hussain_corodet_2021, lee_validation_2023}.
W dziedzinie gastroenterologii wykorzystano CNN do przetwarzania obrazów z endoskopii i ultrasonografii w celu skutecznego wykrywania nieprawidłowości takich jak refluks żołądkowo-przełykowy i zanikowe zapalenie błony śluzowej żołądka \cite{briganti_artificial_2020, yang_application_2019}.
CADx okazało się również skuteczną metodą w przypadku klasyfikacji komórek zarażonych malarią oraz rozpoznawania choroby zwyrodnieniowej stawu kolanowego \cite{rajaraman_pre-trained_2018, al-rimy_adaptive_2023}. 
Liczne badania poświęcone zastosowaniu metod sztucznej inteligencji w diagnostyce medycznej obrazowej wskazują na wysoką skuteczność modeli głębokiego uczenia w zadaniach klasyfikacji binarnej oraz wieloklasowej, wykrywania obiektów na zdjęciach oraz ich segmentacji \cite{ais.2021.15}.
Wynika to z faktu, że systemy wizji komputerowej, w tym konwolucyjne sieci neuronowe, uczą się bezpośrednio z obrazów, czego konsekwencją jest identyfikacja skomplikowanych wzorców i zależności często niezauważanych przez doświadczonych lekarzy.
Modele te funkcjonują jako systemy typu "black-box", co pozwala im zakwalifikować subtelne przypadki chorób i wykryć je w wczesnym stadium \cite{lidstromer_artificial_2021}. 


\section{Computer-aided diagnosis w dermatologii}
Szczególne zastosowanie wizji komputerowej w computer-aided diagnosis stanowi dermatologia, której tematu podjęto się w ramach tej pracy. 


\section{Podsumowanie}
